% ১৪.১২ Loop Statement
\section{Loop Statement}

\begin{learningbox}
এই টপিক শেষে শিক্ষার্থী \texttt{for}, \texttt{while}, \texttt{do...while} loop-এর ব্যবহার লিখতে পারবে।
\end{learningbox}

\begin{definitionbox}{Loop}
একই statement বা statement block বারবার execute করার জন্য যে control structure ব্যবহৃত হয়, তাকে loop বলা হয়।
\end{definitionbox}

\begin{center}
\begin{tabular}{|p{0.20\textwidth}|p{0.44\textwidth}|p{0.24\textwidth}|}
\hline
\textbf{Loop} & \textbf{ধারণা} & \textbf{Condition check} \\
\hline
\texttt{for} & count জানা থাকলে সুবিধাজনক & আগে \\
\hline
\texttt{while} & condition true থাকা পর্যন্ত চলে & আগে \\
\hline
\texttt{do...while} & অন্তত একবার চলে & পরে \\
\hline
\end{tabular}
\end{center}

\begin{verbatim}
for (int i = 1; i <= 5; i++) {
    printf("%d ", i);
}
\end{verbatim}

\begin{mistakebox}{সাধারণ ভুল}
Loop condition update না করলে infinite loop হতে পারে। initialization, condition, increment/decrement খেয়াল রাখতে হবে।
\end{mistakebox}
